\section{Uniform control of the local weights}\label{sec:local}

The ideal sum in Proposition~\ref{prop:unfolding} involves arbitrary
local modules. We now isolate two estimates that are uniform in
those modules and in their indices. The first gives the total
normalization; the second makes large local norm contributions
negligible.

Let $\mathbf k=(\mathbf k_\ell)_{\ell\in S}$ range over nonnegative
valuation tuples, and write
\[
 d_\ell(\mathbf k_\ell)
   =\ell^{\sum_{\mathfrak p\mid\ell}f_{\mathfrak p}k_{\mathfrak p}},
 \qquad d(\mathbf k)=\prod_{\ell\in S}d_\ell(\mathbf k_\ell),
 \qquad P(\mathbf k)=\prod_{\ell\in S}P_\ell(\mathbf k_\ell).
\]
Thus $d$ is the norm of the $S$-supported part of an ideal. Also set
\[
 Z_\ell=\prod_{\mathfrak p\mid\ell}
          (1-\ell^{-f_{\mathfrak p}})^{-1},\qquad
 Z_S=\prod_{\ell\in S}Z_\ell.
\]
Empty products are one.

\begin{lemma}\label{lem:local-moment}
For the weights just defined,
\begin{equation}\label{eq:shell-sum}
 \sum_{\mathbf k}\frac{P(\mathbf k)}{d(\mathbf k)}=\frac{Z_S}{q}.
\end{equation}
If $s=1/(2n)$, then, for every $\eps>0$,
\begin{equation}\label{eq:weight-moment}
 \sum_{\mathbf k}P(\mathbf k)d(\mathbf k)^{s-1}
       \ll_{n,\eps}\frac{q^\eps}{q}.
\end{equation}
In particular,
\begin{equation}\label{eq:weight-tail}
 \sum_{d(\mathbf k)>Q^{1/8}}\frac{P(\mathbf k)}{d(\mathbf k)}
       \ll_n \frac1q Q^{-s/16}.
\end{equation}
\end{lemma}

\begin{proof}
Normalize additive Haar measure on $R_\ell$ to have mass one.
The valuation shell $\mathbf k_\ell$ has mass
$Z_\ell^{-1}/d_\ell(\mathbf k_\ell)$. Its conditional law is
invariant under the transitive unit action, so $P_\ell$ is
exactly the fraction of this shell lying in $L_\ell$.
Since $L_\ell$ has mass $1/q_\ell$, summing over shells and
then multiplying over $S$ proves \eqref{eq:shell-sum}.
Elements with a zero component have additive measure zero.

Choose an additive-Haar uniform element of $L_\ell$, and denote
its component valuations by $v_{\mathfrak p}$. The condition
$R_\ell L_\ell=R_\ell$ implies that the projection to every
component contains a unit. Modulo the $k$th power of its maximal
ideal, the projection therefore contains all $\Z_\ell$ multiples
of that unit, at least $\ell^{\lceil k/e_{\mathfrak p}\rceil}$
elements. Uniform measure on $L_\ell$ induces uniform measure
on this finite additive image. Consequently
\begin{equation}\label{eq:valuation-tail}
 \Pr(v_{\mathfrak p}\ge k)
       \le\ell^{-\lceil k/e_{\mathfrak p}\rceil}\qquad(k\ge1).
\end{equation}
This argument does not require the component projections to be
independent, or even to equal their full integer rings.

Put $j=\max_{\mathfrak p\mid\ell}
\lceil v_{\mathfrak p}/e_{\mathfrak p}\rceil$. Then
$d_\ell(\mathbf v_\ell)\le\ell^{nj}$. If $j\ge a\ge1$, some
$v_{\mathfrak p}$ is at least $e_{\mathfrak p}(a-1)+1$.
Equation~\eqref{eq:valuation-tail} and a union bound give
$\Pr(j\ge a)\le n\ell^{-a}$. Hence
\begin{equation}\label{eq:local-moment}
 \mathbb E_{L_\ell}d_\ell(\mathbf v_\ell)^s
 \le1+n\sum_{a\ge1}\ell^{(ns-1)a}
 \le1+C_n\ell^{-1/2}.
\end{equation}
Using the shell masses once more yields the exact factorization
\[
 \sum_{\mathbf k}P(\mathbf k)d(\mathbf k)^{s-1}
   =\frac1q\prod_{\ell\in S}
       \left(Z_\ell\,\mathbb E_{L_\ell}d_\ell(\mathbf v_\ell)^s\right).
\]
Each factor is at most
$(1-\ell^{-1})^{-n}(1+C_n\ell^{-1/2})$. For fixed $\eps>0$
this is at most $\ell^\eps$ at all sufficiently large primes;
the finitely many remaining primes contribute a constant
depending only on $n,\eps$. As $\prod_{\ell\in S}\ell\le q$,
this proves \eqref{eq:weight-moment}. Finally, take
$\eps=s/16$ and use $q\le Q$ to obtain
\[
 \sum_{d>Q^{1/8}}\frac{P(\mathbf k)}d
 \le Q^{-s/8}\sum_{\mathbf k}P(\mathbf k)d^{s-1}
 \ll_n q^{-1}Q^{-s/16}.
 \qedhere
\]
\end{proof}
